\clearpage

\section{ANTECEDENTES}

El Desafío Bebras cuenta con antecedentes académicos e institucionales que lo posicionan como una referencia internacional en la promoción del pensamiento computacional en estudiantes de nivel escolar. Uno de los trabajos fundacionales es el de Dagien\.e y Futschek, quienes describen a Bebras como una competencia orientada a acercar conceptos fundamentales de la informática mediante tareas breves, atractivas y resolubles sin conocimientos previos de programación, además de proponer criterios para el diseño de tareas de calidad dentro del desafío \cite{dagienefutschek2008}. A nivel organizativo, la comunidad Bebras consolidó posteriormente un modelo internacional de colaboración basado en talleres anuales de discusión y refinamiento de tareas, lo que permitió mantener lineamientos comunes y, al mismo tiempo, facilitar que cada país adapte el concurso a su propio contexto educativo \cite{bebras,bebrastaskworkshops,bebrasstructure}.

En el plano educativo, distintos estudios ampliaron estos antecedentes al analizar el valor pedagógico de Bebras más allá de la competencia anual. Dagien\.e y Stupurien\.e explican que el modelo evolucionó de un concurso puntual a una comunidad educativa sostenible orientada al aprendizaje de conceptos de informática y pensamiento computacional \cite{dagienestupuriene2016}. De forma complementaria, Izu et al. realizaron un estudio multinacional sobre el contenido y el rendimiento de las tareas Bebras, mostrando que estas permiten examinar categorías de pensamiento computacional y desempeño estudiantil a gran escala \cite{izumirolo2017}. En conjunto, estos trabajos aportan una base conceptual y metodológica sólida para considerar a Bebras como referencia válida en el diseño de propuestas educativas apoyadas en tecnología.

Desde el punto de vista tecnológico, los antecedentes muestran que Bebras no funciona mediante una única plataforma centralizada, sino a través de sistemas nacionales de gestión del concurso. La propia organización internacional documenta la existencia de distintos \textit{contest management systems} para la administración del desafío, mientras que experiencias como \textit{Castor Informatique} en Francia, \textit{Bebras Perú} y \textit{Bebras Computing Challenge} en Estados Unidos evidencian que los países desarrollan o adaptan plataformas web para el registro de participantes, la administración de tareas, la práctica previa y la ejecución del concurso \cite{bebrascms,bebrasfrancia,bebrasperu,bebraseeuu}.